% !TeX document-id = {58dd56a7-a545-4332-9f31-aa9515ad78e0}
%%%%%%%%%%%%%%%%%%%%%%%%%%%%%%%%%%%%%%%%%
% The Legrand Orange Book
% LaTeX Template
% Version 2.4 (26/09/2018)
%
% This template was downloaded from:
% http://www.LaTeXTemplates.com
%
% Original author:
% Mathias Legrand (legrand.mathias@gmail.com) with modifications by:
% Vel (vel@latextemplates.com)
%
% License:
% CC BY-NC-SA 3.0 (http://creativecommons.org/licenses/by-nc-sa/3.0/)
%
% Compiling this template:
% This template uses biber for its bibliography and makeindex for its index.
% When you first open the template, compile it from the command line with the 
% commands below to make sure your LaTeX distribution is configured correctly:
%
% 1) pdflatex main
% 2) makeindex main.idx -s StyleInd.ist
% 3) biber main
% 4) pdflatex main x 2
%
% After this, when you wish to update the bibliography/index use the appropriate
% command above and make sure to compile with pdflatex several times 
% afterwards to propagate your changes to the document.
%
% This template also uses a number of packages which may need to be
% updated to the newest versions for the template to compile. It is strongly
% recommended you update your LaTeX distribution if you have any
% compilation errors.
%
% Important note:
% Chapter heading images should have a 2:1 width:height ratio,
% e.g. 920px width and 460px height.
%
%%%%%%%%%%%%%%%%%%%%%%%%%%%%%%%%%%%%%%%%%

%----------------------------------------------------------------------------------------
%	PACKAGES AND OTHER DOCUMENT CONFIGURATIONS
%----------------------------------------------------------------------------------------

% !BIB program = biber
% !TeX TXS-program:compile = txs:///pdflatex/[--shell-escape]

\documentclass[11pt,fleqn]{book} % Default font size and left-justified equations

\input{structure.tex} % Insert the commands.tex file which contains the majority of the structure behind the template

\hypersetup{pdftitle={Arduboy Starters Guide},pdfauthor={Stijn Caerts}} % Uncomment and fill out to include PDF metadata for the author and title of the book

\usepackage[scale=2]{ccicons}

\def\Cpp{{C\nolinebreak[4]\hspace{-.05em}\raisebox{.4ex}{\tiny\bf ++}}}

%----------------------------------------------------------------------------------------

\begin{document}

%----------------------------------------------------------------------------------------
%	TITLE PAGE
%----------------------------------------------------------------------------------------

\begingroup
\thispagestyle{empty} % Suppress headers and footers on the title page
\begin{tikzpicture}[remember picture,overlay]
\node[inner sep=0pt] (background) at (current page.center) {\includegraphics[width=\paperwidth]{jcw_background.pdf}};
\draw (current page.center) node [fill=ocre!30!white,fill opacity=0.9,text opacity=1,inner sep=1cm]{\Huge\centering\bfseries\sffamily\parbox[c][][t]{\paperwidth}{\centering Arduboy\\[15pt] % Book title
{\Large Starters Guide}\\[20pt] % Subtitle
{\huge Stijn Caerts}}}; % Author name

\node[yshift=2cm] at (current page.south){\includegraphics[height=3cm]{jcw_logo.pdf}};
\end{tikzpicture}
\vfill
\endgroup

%----------------------------------------------------------------------------------------
%	COPYRIGHT PAGE
%----------------------------------------------------------------------------------------

\newpage
~\vfill
\thispagestyle{empty}
{
\fontfamily{ptm}\selectfont

\noindent Copyright \copyright\ 2019 Stijn Caerts\\ % Copyright notice

\noindent \textsc{Jeugd, Cultuur en Wetenschap vzw}\\ % Publisher

\noindent \textsc{\href{https://stijn.caerts.be/}{stijn.caerts.be}} -- \textsc{\href{https://www.jeugdcultuurenwetenschap.be/}{www.jeugdcultuurenwetenschap.be}}\\ % URL

\noindent Dit werk valt onder een \href{https://creativecommons.org/licenses/by-nc-sa/4.0/deed.nl}{Creative Commons Naamsvermelding-NietCommercieel-GelijkDelen 4.0 Internationaal-licentie} (de ``Licentie''). Dit document mag enkel gebruikt worden in navolging van de Licentie. De volledige Licentie-tekst is beschikbaar op \url{https://creativecommons.org/licenses/by-nc-sa/4.0/}.\\

\noindent \href{https://creativecommons.org/licenses/by-nc-sa/4.0/}{\ccbyncsaeu}

\vspace{1cm}

\noindent \textit{Eerste versie, augustus 2019} % Printing/edition date
}

%----------------------------------------------------------------------------------------
%	TABLE OF CONTENTS
%----------------------------------------------------------------------------------------

%\usechapterimagefalse % If you don't want to include a chapter image, use this to toggle images off - it can be enabled later with \usechapterimagetrue

\chapterimage{jcw_header.pdf} % Table of contents heading image

\pagestyle{empty} % Disable headers and footers for the following pages

\tableofcontents % Print the table of contents itself

\cleardoublepage % Forces the first chapter to start on an odd page so it's on the right side of the book

\pagestyle{fancy} % Enable headers and footers again

%----------------------------------------------------------------------------------------
%	PART
%----------------------------------------------------------------------------------------

\part{Introductie tot C++}
\chapterimage{jcw_header.pdf}
\chapter{Inleiding}
\section{Wat is C++?}
% https://www.tutorialspoint.com/cplusplus
Programma's voor \textbf{Arduino} en \textbf{Arduboy} worden geschreven in de programmeertaal \Cpp\index{C++@\Cpp}. Het is niet nodig om de hele programmeertaal te kennen en begrijpen voor je aan de slag kan gaan met programmeren. Daarom geven we hier een beknopt overzicht van de belangrijkste concepten die je nodig hebt om van start te gaan.\\

\noindent In de volgende hoofdstukken komen variabelen en types, controlestructuren (if-then-else, for, while) en functies en procedures aan bod. Tot slot zijn er nog twee hoofdstukken die dieper ingaan op de mogelijkheden van \Cpp, namelijk arrays en lijsten, en klassen en objecten.\\

\noindent Voor een interactieve en uitgebreidere introductie tot \Cpp, kan je terecht bij W3Schools (\url{https://www.w3schools.com/cpp/}).

% TODO compiler

\section{Syntax}
\subsection{Puntkomma}
Achter elke instructie wordt in \Cpp{} een puntkomma \mintinline{cpp}{;} geplaatst. Deze puntkomma vertelt de compiler dat op die plaats een instructie eindigt. Als je een puntkomma vergeet te plaatsen, is het programma niet correct en zal je het programma niet kunnen compileren. De compiler zal dan een foutmelding geven.

\subsection{Accolades}
Accolades \mintinline{cpp}{{}} worden gebruikt om een blok code aan te duiden. De accolades moeten gebalanceerd voorkomen in de code, wat betekent dat voor elke openende accolade \texttt{\{} er een bijhorende sluitende accolade \texttt{\}} moet zijn. Het gebruik van accolades wordt later geïllustreerd bij de controlestructuren (if-then-else, for, ..), functies en klassen.

\subsection{Commentaar}
% commentaar in C++
% // of /* */
Om de leesbaarheid van je code te verhogen, is het nuttig om commentaar toe te voegen. In deze commentaar\index{Commentaar} beschrijf je wat dit deel van de code juist doet. Hierdoor is het duidelijk wat je juist hebt geprogrammeerd, ook als je later opnieuw je code bekijkt.\\

\noindent In \Cpp{} zijn er twee verschillende manieren om commentaar toe te voegen.\\
Commentaar op één lijn wordt aangeduid met \mintinline{cpp}{//}. Alle tekst na \mintinline{cpp}{//} tot het einde van de lijn wordt beschouwd als commentaar en zal bijgevolg niet uitgevoerd worden. 

\begin{example}[Commentaarlijn]
	\phantom{ }
	\begin{minted}{cpp}
// Dit is een lijn commentaar
int a = 42;
	\end{minted}
\end{example}

\noindent Commentaar over meerdere lijnen start met \texttt{/*} en eindigt met \texttt{*/}. Alle tekst tussen \texttt{/*} en \texttt{*/} wordt door de compiler genegeerd.

\begin{example}[Commentaar over meerdere lijnen]
	\phantom{ }
	\begin{minted}{cpp}
/*
  Deze commentaar neemt
  meerdere lijnen in beslag.
*/
int a = 42;
	\end{minted}
\end{example}

\chapter{Variabelen en types}
\index{Variabele}
\index{Datatype}
\index{Type|seealso{Datatype}}
Variabelen worden gebruikt om informatie of gegevens bij te houden in het programma. Een variabele is een plaats in het geheugen, met een zelfgekozen naam, waar gegevens in opgeslagen kunnen worden. In \Cpp{} heeft elke variabele een type. Een type komt overeen met het soort van gegevens dat opgeslagen kunnen worden in zo'n variabele.

\section{Datatypes}
% https://www.tutorialspoint.com/cplusplus/cpp_data_types.htm
% strings https://www.tutorialspoint.com/cplusplus/cpp_strings

Op basis van het datatype van een variabele wordt er geheugen gereserveerd om informatie in op te slaan. In \Cpp{} zijn de volgende primitieve datatypes gedefinieerd:

\begin{definition}[Primitieve datatypes]
	\phantom{ } \\
	\begin{minipage}{\columnwidth}
		\vspace{0.1cm}
		\begin{tabular}{rl}
			\textbf{bool} & Booleaanse waarde: \texttt{true} of \texttt{false} \\
			\textbf{char} & Karakter \\
			\textbf{int} & Geheel getal \\
			\textbf{float} & Kommagetal (met enkele precisie) \\
			\textbf{double} & Kommagetal (met dubbele precisie)
		\end{tabular}
		\vspace{0.1cm}
	\end{minipage}
\end{definition}

\index{Datatype!bool}
\index{Datatype!char}
\index{Datatype!int}
\index{Datatype!float}
\index{Datatype!double}

\index{Booleaanse waarde|seealso{Datatype, bool}}

\noindent Deze types worden gebruikt bij variabelen, maar ook bij functies. Deze primitieve datatypes kunnen ook gebruikt worden om complexere types samen te stellen, zoals bijvoorbeeld bij arrays en klassen.

\section{Variabelen}
Een variabele kan gezien worden als een opslagplaats waaraan een naam verbonden is. In deze opslagplaats kan een waarde geplaatst worden.
\subsection{Declareren}
\index{Variabele!Declareren}
% type naam_variabele;
Alvorens informatie kan opgeslagen worden in een variabele, moet de variabele eerst \emph{gedefinieerd} of \emph{gedeclareerd} worden. De declaratie van een variabele vertelt de compiler dat er plaats in het geheugen gereserveerd moet worden om gegevens van het gegeven datatype in op te slaan.\\

\noindent Een nieuwe variabele declareren doe je met een instructie van deze vorm: 
\begin{center}
	\texttt{type naam\_variabele;}
\end{center}

\begin{example}[Declareren van een \emph{integer} variabele]
	\phantom{ }
	\begin{minted}{cpp}
// Declareer een nieuwe variabele van het type int met naam a
int a;
	\end{minted}
\end{example}

\noindent Je kan bij het definiëren van een nieuwe variabele er ook meteen een waarde aan toekennen.

\begin{example}[Declareren van een \emph{integer} variabele en een waarde toekennen]
	\phantom{ }
	\begin{minted}{cpp}
/*
  Declareer een nieuwe variabele van het type int met naam a, 
  en ken de waarde 42 toe
*/
int a = 42;
	\end{minted}
\end{example}

\noindent Bij het kiezen van een variabelenaam moet je met de volgende regels rekening houden:
\begin{itemize}
	\item namen moeten beginnen met een letter of een laag streepje (\_);
	\item namen kunnen letters, cijfers en lage streepjes bevatten;
	\item namen zijn hoofdlettergevoelig;
	\item spaties of speciale karakters zijn niet toegelaten;
	\item \Cpp{}-keywords (zoals \texttt{int}, \texttt{if}, ...) kunnen niet gebruikt worden als variabelenaam.
\end{itemize}

\todo{Voorbeeld van elk type + waarde}

\subsection{Scope}
\index{Variabele!Scope}
De \emph{scope} van een variabele is de plaats in het programma waar je een gedefinieerde variabele kan gebruiken. Je kan dit ook beschouwen als de levensduur van de variabele.\\

\noindent
De scope van een variabele wordt bepaald door waar de variabele werd gedeclareerd. Dit kunnen we opdelen in 3 verschillende soorten plaatsen.

\begin{enumerate}
	\item in een functie of blok: \textbf{lokale variabelen},
	\item in de definitie van een functie: \textbf{parameters},
	\item buiten alle functies: \textbf{globale variabelen}.
\end{enumerate}

\noindent
In het algemeen kan gesteld worden dat de scope van de variabelen beperkt is tot het blok, aangeduid met accolades \texttt{\{\}}, waarin de variabele gedefinieerd is.

\begin{example}[Scope van variabelen]
	\phantom{ }
	\begin{minted}{cpp}
// Globale variabele g
int g = 12;

void setup() {
  // Lokale variabelen a en b
  int a = 42;
  
  
  /*
    Globale variabelen kunnen na declaratie 
    overal gebruikt worden
  */
  int b = a + g;
  
  // Na declaratie moet het type van variabelen niet gespecifieerd worden
  g = b;
}

// Variabelen a en b kunnen buiten de functie setup() niet gebruikt worden

	\end{minted}
\end{example}

\subsection{Operatoren}
Om de waarde in de variabelen aan te passen, maken we gebruik van \emph{operatoren}. Hieronder geven we een overzicht van de belangrijkste operatoren in \Cpp.

\index{Operatoren}
\subsubsection{Rekenkundige operatoren}
\index{Operatoren!Rekenkundig}

Rekenkundige operatoren worden gebruikt om veelgebruikte wiskundige operaties uit te voeren.
In deze voorbeelden gebruiken we de variabelen $ A = 10 $ en $ B = 3 $.

\begin{center}
	\begin{tabular}{c l l}
		\toprule
		\textbf{Operator} & \textbf{Beschrijving}                     & \textbf{Voorbeeld}                                              \\ \midrule
		   \texttt{+}     & Optelling                                 & \mintinline{cpp}{A + B} $\rightarrow 13$                        \\
		   \texttt{-}     & Vermindering                              & \mintinline{cpp}{A - B} $\rightarrow 7$                         \\
		   \texttt{*}     & Vermenigvuldiging                         & \mintinline{cpp}{A * B} $\rightarrow 30$                        \\
		   \texttt{/}     & Deling                                    & \mintinline{cpp}{A / B} $\rightarrow 2$                         \\
		   \texttt{\%}    & Modulus (rest na deling)                  & \mintinline{cpp}{A % B} $\rightarrow 1$                         \\
		   \texttt{++}    & Verhoogt de waarde van de variabele met 1 & \mintinline{cpp}{A++} $\Rightarrow$ \texttt{A} $\rightarrow 11$ \\
		  \texttt{-{}-}   & Verlaagt de waarde van de variabele met 1 & \mintinline{cpp}{A--} $\Rightarrow$ \texttt{A} $\rightarrow 9$  \\ \bottomrule
	\end{tabular}
\end{center}


\subsubsection{Toewijzende operatoren}
\index{Operatoren!Toewijzend}

De operator voor de toewijzing van een waarde aan een variabele is het gelijkheidsteken (\texttt{=}). Naast deze toewijzingsoperator bestaan er nog andere varianten, die kortere notaties zijn voor een rekenkundige operatie gevolgd door een toewijzing.

\begin{center}
	\begin{tabular}{c c l}
		\toprule
		\textbf{Operator} & \textbf{Voorbeeld}                     & \textbf{Equivalent}           \\ \midrule
		\texttt{=} & \mintinline{cpp}{A = 12} & \mintinline{cpp}{A = 12}\\
		\texttt{+=} & \mintinline{cpp}{A += 2} & \mintinline{cpp}{A = A + 2} \\
		\texttt{-=} & \mintinline{cpp}{A -= 3} & \mintinline{cpp}{A = A - 3} \\
		\texttt{*=} & \mintinline{cpp}{A *= 4} & \mintinline{cpp}{A = A * 4} \\
		\bottomrule
	\end{tabular}
\end{center}

\subsubsection{Vergelijkende operatoren}
\index{Operatoren!Vergelijkend}
Vergelijkende operatoren worden gebruikt om twee waardes met elkaar te vergelijken. Het resultaat van deze operatie is een Booleaanse waarde, \texttt{true} of \texttt{false}.

\begin{center}
	\begin{tabular}{c l l}
		\toprule
		\textbf{Operator} & \textbf{Beschrijving}     & \textbf{Voorbeeld}                                     \\ \midrule
		   \texttt{==}    & Gelijk aan                & \mintinline{cpp}{10 == 3} $\rightarrow$ \texttt{false} \\
		   \texttt{!=}    & Niet gelijk aan           & \mintinline{cpp}{10 != 3} $\rightarrow$ \texttt{true}  \\
		   \texttt{>}     & Groter dan                & \mintinline{cpp}{10 > 10} $\rightarrow$ \texttt{false} \\
		   \texttt{<}     & Kleiner dan               & \mintinline{cpp}{5 < 10} $\rightarrow$ \texttt{true}   \\
		   \texttt{>=}    & Groter dan of gelijk aan  & \mintinline{cpp}{10 >= 10} $\rightarrow$ \texttt{true} \\
		   \texttt{<=}    & Kleiner dan of gelijk aan & \mintinline{cpp}{5 <= 10} $\rightarrow$ \texttt{true}  \\ \bottomrule
	\end{tabular}
\end{center}

\subsubsection{Logische operatoren}
\index{Operatoren!Logisch}
Logische operatoren worden gebruikt om meerdere Booleaanse waarden te combineren in een formule tot één waarheidswaarde.

\begin{center}
	\begin{tabular}{c l l}
		\toprule
		\textbf{Operator} & \textbf{Beschrijving} & \textbf{Voorbeeld} \\ \midrule
		\texttt{\&\&} & Logische conjunctie (\textbf{AND}) & \mintinline{cpp}{(true && false)} $\rightarrow$ \texttt{false} \\
		\texttt{||} & Logische disjunctie (\textbf{OR}) & \mintinline{cpp}{(true || false)} $\rightarrow$ \texttt{true} \\
		\texttt{!} & Logische negatie (\textbf{NOT}) & \mintinline{cpp}{! true} $\rightarrow$ \texttt{false} \\
		\bottomrule
	\end{tabular}
\end{center}

\noindent Voor meer voorbeelden van het gebruik van logische operatoren, kan je een kijkje nemen naar waarheidstabellen~\cite{wiki:Waarheidstabel}.


\chapter{Controlestructuren}
Controlestructuren bepalen de loop van het programma. Zo is het mogelijk om stukken code enkel in specifieke gevallen uit te voeren, of om een blok code meerdere keren te herhalen.
% if-then-else, switch, for, while, (break, continue)
\section{If-then-else}
De eenvoudigste controlestructuur is het \emph{if}-statement. Dit kan verder uitgebreid worden met een \emph{else}-blok voor verdere controle over de uitvoeringsvolgorde.

\subsection{If-then}

\begin{definition}[If-statement]
	\phantom{ }
	\begin{minted}{cpp}
  if (conditie) {
    /*
      instructies tussen deze accolades worden enkel 
      uitgevoerd als de conditie tot true evalueert
    */
  }
	\end{minted}
\end{definition}

\subsection{Else}
Na een if-statement kan optioneel een else-blok geplaatst worden. De instructies in dit blok zullen enkel uitgevoerd worden als aan de conditie van het if-statement \textbf{niet} voldaan is. De code ziet er dan uit als volgt.

\begin{definition}[If-then-else]
	\phantom{ }
	\begin{minted}{cpp}
  if (conditie) {
    // conditie evalueert tot true
  } else {
    // conditie evalueert tot false
  }
	\end{minted}
\end{definition}

\subsection{Geneste if-statements}

\section{Switch}

\section{While-lus}

\section{For-lus}


\chapter{Functies en procedures}
% https://www.dinkum.nl/software/cplusplus/doc/functions/
\chapter{Arrays en lijsten \textsc{\footnotesize(Advanced)}}
% https://www.dinkum.nl/software/cplusplus/doc/arrays/
% https://github.com/ivanseidel/LinkedList
\chapter{Klassen en objecten \textsc{\footnotesize(Advanced)}}
% https://www.dinkum.nl/software/cplusplus/doc/OO/classes.html
% Geheugenbeheer
% https://stackoverflow.com/questions/6337294/creating-an-object-with-or-without-new

\part{Arduboy}

\chapter{Arduino}
\section{Programmastructuur}
\subsection{Globale variabelen}
\subsection{De \texttt{setup()} procedure}
\subsection{De \texttt{loop()} procedure}

\section{Importeren van libraries}

\chapter{Arduboy}
\section{Instellingen}
% Arduino Leonardo
% Importeren library

\section{De Arduboy2 library}
\subsection{Display}
\subsection{Buttons}

\section{Emulator}
\section{Programma op Arduboy plaatsen}



\part{Part One}

%----------------------------------------------------------------------------------------
%	CHAPTER 1
%----------------------------------------------------------------------------------------

\chapter{Text Chapter}

\section{Paragraphs of Text}\index{Paragraphs of Text}

\lipsum[1-7] % Dummy text

%------------------------------------------------

\section{Citation}\index{Citation}

This statement requires citation \cite{article_key}; this one is more specific \cite[162]{book_key}.

%------------------------------------------------

\section{Lists}\index{Lists}

Lists are useful to present information in a concise and/or ordered way\footnote{Footnote example...}.

\subsection{Numbered List}\index{Lists!Numbered List}

\begin{enumerate}
\item The first item
\item The second item
\item The third item
\end{enumerate}

\subsection{Bullet Points}\index{Lists!Bullet Points}

\begin{itemize}
\item The first item
\item The second item
\item The third item
\end{itemize}

\subsection{Descriptions and Definitions}\index{Lists!Descriptions and Definitions}

\begin{description}
\item[Name] Description
\item[Word] Definition
\item[Comment] Elaboration
\end{description}

%----------------------------------------------------------------------------------------
%	CHAPTER 2
%----------------------------------------------------------------------------------------

\chapter{In-text Elements}

\section{Theorems}\index{Theorems}

This is an example of theorems.

\subsection{Several equations}\index{Theorems!Several Equations}
This is a theorem consisting of several equations.

\begin{theorem}[Name of the theorem]
In $E=\mathbb{R}^n$ all norms are equivalent. It has the properties:
\begin{align}
& \big| ||\mathbf{x}|| - ||\mathbf{y}|| \big|\leq || \mathbf{x}- \mathbf{y}||\\
&  ||\sum_{i=1}^n\mathbf{x}_i||\leq \sum_{i=1}^n||\mathbf{x}_i||\quad\text{where $n$ is a finite integer}
\end{align}
\end{theorem}

\subsection{Single Line}\index{Theorems!Single Line}
This is a theorem consisting of just one line.

\begin{theorem}
A set $\mathcal{D}(G)$ in dense in $L^2(G)$, $|\cdot|_0$. 
\end{theorem}

%------------------------------------------------

\section{Definitions}\index{Definitions}

This is an example of a definition. A definition could be mathematical or it could define a concept.

\begin{definition}[Definition name]
Given a vector space $E$, a norm on $E$ is an application, denoted $||\cdot||$, $E$ in $\mathbb{R}^+=[0,+\infty[$ such that:
\begin{align}
& ||\mathbf{x}||=0\ \Rightarrow\ \mathbf{x}=\mathbf{0}\\
& ||\lambda \mathbf{x}||=|\lambda|\cdot ||\mathbf{x}||\\
& ||\mathbf{x}+\mathbf{y}||\leq ||\mathbf{x}||+||\mathbf{y}||
\end{align}
\end{definition}

%------------------------------------------------

\section{Notations}\index{Notations}

\begin{notation}
Given an open subset $G$ of $\mathbb{R}^n$, the set of functions $\varphi$ are:
\begin{enumerate}
\item Bounded support $G$;
\item Infinitely differentiable;
\end{enumerate}
a vector space is denoted by $\mathcal{D}(G)$. 
\end{notation}

%------------------------------------------------

\section{Remarks}\index{Remarks}

This is an example of a remark.

\begin{remark}
The concepts presented here are now in conventional employment in mathematics. Vector spaces are taken over the field $\mathbb{K}=\mathbb{R}$, however, established properties are easily extended to $\mathbb{K}=\mathbb{C}$.
\end{remark}

%------------------------------------------------

\section{Corollaries}\index{Corollaries}

This is an example of a corollary.

\begin{corollary}[Corollary name]
The concepts presented here are now in conventional employment in mathematics. Vector spaces are taken over the field $\mathbb{K}=\mathbb{R}$, however, established properties are easily extended to $\mathbb{K}=\mathbb{C}$.
\end{corollary}

%------------------------------------------------

\section{Propositions}\index{Propositions}

This is an example of propositions.

\subsection{Several equations}\index{Propositions!Several Equations}

\begin{proposition}[Proposition name]
It has the properties:
\begin{align}
& \big| ||\mathbf{x}|| - ||\mathbf{y}|| \big|\leq || \mathbf{x}- \mathbf{y}||\\
&  ||\sum_{i=1}^n\mathbf{x}_i||\leq \sum_{i=1}^n||\mathbf{x}_i||\quad\text{where $n$ is a finite integer}
\end{align}
\end{proposition}

\subsection{Single Line}\index{Propositions!Single Line}

\begin{proposition} 
Let $f,g\in L^2(G)$; if $\forall \varphi\in\mathcal{D}(G)$, $(f,\varphi)_0=(g,\varphi)_0$ then $f = g$. 
\end{proposition}

%------------------------------------------------

\section{Examples}\index{Examples}

This is an example of examples.

\subsection{Equation and Text}\index{Examples!Equation and Text}

\begin{example}
Let $G=\{x\in\mathbb{R}^2:|x|<3\}$ and denoted by: $x^0=(1,1)$; consider the function:
\begin{equation}
f(x)=\left\{\begin{aligned} & \mathrm{e}^{|x|} & & \text{si $|x-x^0|\leq 1/2$}\\
& 0 & & \text{si $|x-x^0|> 1/2$}\end{aligned}\right.
\end{equation}
The function $f$ has bounded support, we can take $A=\{x\in\mathbb{R}^2:|x-x^0|\leq 1/2+\epsilon\}$ for all $\epsilon\in\intoo{0}{5/2-\sqrt{2}}$.
\end{example}

\subsection{Paragraph of Text}\index{Examples!Paragraph of Text}

\begin{example}[Example name]
\lipsum[2]
\end{example}

%------------------------------------------------

\section{Exercises}\index{Exercises}

This is an example of an exercise.

\begin{exercise}
This is a good place to ask a question to test learning progress or further cement ideas into students' minds.
\end{exercise}

%------------------------------------------------

\section{Problems}\index{Problems}

\begin{problem}
What is the average airspeed velocity of an unladen swallow?
\end{problem}

%------------------------------------------------

\section{Vocabulary}\index{Vocabulary}

Define a word to improve a students' vocabulary.

\begin{vocabulary}[Word]
Definition of word.
\end{vocabulary}

%----------------------------------------------------------------------------------------
%	PART
%----------------------------------------------------------------------------------------

\part{Part Two}

%----------------------------------------------------------------------------------------
%	CHAPTER 3
%----------------------------------------------------------------------------------------

\chapter{Presenting Information}

\section{Table}\index{Table}

\begin{table}[h]
\centering
\begin{tabular}{l l l}
\toprule
\textbf{Treatments} & \textbf{Response 1} & \textbf{Response 2}\\
\midrule
Treatment 1 & 0.0003262 & 0.562 \\
Treatment 2 & 0.0015681 & 0.910 \\
Treatment 3 & 0.0009271 & 0.296 \\
\bottomrule
\end{tabular}
\caption{Table caption}
\label{tab:example} % Unique label used for referencing the table in-text
%\addcontentsline{toc}{table}{Table \ref{tab:example}} % Uncomment to add the table to the table of contents
\end{table}

Referencing Table \ref{tab:example} in-text automatically.

%------------------------------------------------

\section{Figure}\index{Figure}

\begin{figure}[h]
\centering\includegraphics[scale=0.5]{placeholder.jpg}
\caption{Figure caption}
\label{fig:placeholder} % Unique label used for referencing the figure in-text
%\addcontentsline{toc}{figure}{Figure \ref{fig:placeholder}} % Uncomment to add the figure to the table of contents
\end{figure}

Referencing Figure \ref{fig:placeholder} in-text automatically.

%----------------------------------------------------------------------------------------
%	BIBLIOGRAPHY
%----------------------------------------------------------------------------------------

\chapterimage{jcw_header.pdf}

\chapter*{Bibliografie}
\addcontentsline{toc}{chapter}{\textcolor{ocre}{Bibliografie}} % Add a Bibliography heading to the table of contents

%------------------------------------------------

\section*{Artikels}
\addcontentsline{toc}{section}{Artikels}
\printbibliography[heading=bibempty,type=article]

%------------------------------------------------

\section*{Boeken}
\addcontentsline{toc}{section}{Boeken}
\printbibliography[heading=bibempty,type=book]

\section*{Websites}
\printbibliography[heading=bibempty,type=misc]

%----------------------------------------------------------------------------------------
%	INDEX
%----------------------------------------------------------------------------------------

\cleardoublepage % Make sure the index starts on an odd (right side) page
\phantomsection
\setlength{\columnsep}{0.75cm} % Space between the 2 columns of the index
\addcontentsline{toc}{chapter}{\textcolor{ocre}{Index}} % Add an Index heading to the table of contents
\printindex % Output the index

%----------------------------------------------------------------------------------------

\end{document}
